\section{Distributional localization with network-dependent exceptions}
\label{sec:localized-prior-theorem}

Stability need not hold on the full cube. The exceptional region can even
vary with the network, provided its input-set complexity is controlled.
The prior below is an exactly conditioned Gaussian, rather than a mixture
of separately normalized Gaussian components. Throughout these appendices,
\emph{Gibbs test-margin error} means the posterior-averaged test loss in
\eqref{eq:standalone-risk}.

Let $\mathcal B=\{B_1,\ldots,B_M\}$ be a family of distinct subsets of
$\mathcal X$, containing the empty set. The family may be chosen using an
independent unlabeled pilot. Draw a further unlabeled validation sample
$V$ of size $m$, independent of the labeled sample $S$. Fix $0\le\rho<1$.
Condition the raw Gaussian on the existence of an index $j$ such that
\begin{equation}\label{eq:localized-support}
 \widehat D_V(B_j)\le\rho,\qquad
 |T(z_w(x;\eta))-T(z_w(x))|\le1/16
 \quad(x\notin B_j,\ \|\eta\|_*\le r_c(w)).
\end{equation}
Then apply the unchanged training likelihood $f(X_i)z_w(X_i)>1/2$ at every
training example, including examples in the exceptional region.

\begin{theorem}[Localized stability under exact union conditioning]
\label{thm:localized-union}
For targets covered by Theorem~\ref{thm:standalone-clipped}, the posterior
with support~\eqref{eq:localized-support} has Gibbs test-margin error at
most $\rho+\varepsilon$, with probability at least $1-\delta$, whenever
\[
 m\ge O\!\left(\frac{\log M+\log(1/\delta)}{\varepsilon^2}\right),
 \qquad
 n\ge\operatorname{poly}
   (d,L,C,c+c^{-1},\log M,\varepsilon^{-1},\log(1/\delta)).
\]
One sample event gives the guarantee simultaneously for all $N\ge C$ and
$0<t\le1$. No modification of the training likelihood is required.
\end{theorem}

\begin{proof}
Write $a=N/t^2\ge1$. By a one-sided Hoeffding bound and a union bound,
with probability at least $1-\delta/2$ over $V$,
\[
 D(B_j)\le\widehat D_V(B_j)+u\qquad(1\le j\le M),
\]
provided $m\ge[\log M+\log(2/\delta)]/(2u^2)$.
The global target cylinder is admissible with the empty exceptional set,
so for every $V,S$ the full posterior $Q$ satisfies
$D_{\rm KL}(Q\Vert\mu_{N,t})\le aD_*$.

Assign to each network the first index witnessing
\eqref{eq:localized-support}. Write $q_j$ for its posterior probability
and $Q_j$ for the conditional network law. The assignment is measurable
because the family is finite. The relative-entropy chain rule gives
\begin{equation}\label{eq:localized-component-entropy}
 \sum_jq_jD_{\rm KL}(Q_j\Vert\mu_{N,t})
 =D_{\rm KL}(Q\Vert\mu_{N,t})+H(q)
 \le aD_*+\log M.
\end{equation}
Thus indices for which $D_{\rm KL}(Q_j\Vert\mu_{N,t})>aK_1$ have total
posterior mass at most $\beta$, where
\[
 K_1=\frac{D_*+\log M}{\beta}.
\]
This step avoids assuming that each assigned component separately contains
the target cylinder; that assumption would be false in general.

For each retained component, the Gaussian entropy moment bound implies
$\mathbb E_{Q_j}E_v\le4K_1+2\log2$. Set
$B=(4K_1+2\log2)/\beta$ and discard its readout tail $E_v>B$, of mass at
most $\beta$. Denote the resulting conditional law by $Q'_j$.
For any conditioning event $E$ of probability $p$, KL decomposition gives
\[
 pD_{\rm KL}(Q_j(\cdot\mid E)\Vert\mu_{N,t})
 \le D_{\rm KL}(Q_j\Vert\mu_{N,t})+\log2.
\]
Consequently, for $\beta\le1/2$,
$D_{\rm KL}(Q'_j\Vert\mu_{N,t})\le2a(K_1+1)$.
All retained components are stable outside their own $B_j$ at the common
radius $r=\min\{1/6,c/\sqrt{1+HB}\}$.

Compress each $Q'_j$ separately, using the common finite class from
Lemma~\ref{lem:raw-entropy-compression} with $K=2(K_1+1)$.
Apply the finite-class relative-Chernoff bound simultaneously to the
$M|\mathcal H|$ masked ramp losses
\[
 \ell_{j,h}(x)=\mathbf1\{x\notin B_j\}\,
                      \psi(f(x)T(h(x))).
\]
This event has probability at least $1-\delta/2$ over the labeled sample.
Outside $B_j$, training margin consistency and the compression coupling
make the candidate's empirical masked loss at most $\alpha$ after
averaging. The population comparison adds another $\alpha$.
Therefore every retained component has masked test-margin risk at most
\[
 3\alpha+\frac{8[\log M+\log|\mathcal H|+\log(2/\delta)]}{n}.
\]
On the validation event, its unmasked risk costs at most another
$D(B_j)\le\rho+u$. Averaging the components and restoring both discarded
masses gives
\[
 R(Q)\le\rho+u+2\beta+3\alpha+
       \frac{8[\log M+\log|\mathcal H|+\log(2/\delta)]}{n}.
\]
Take $u=\varepsilon/4$, $\beta=\varepsilon/8$,
$\alpha=\varepsilon/12$, and make the last term at most
$\varepsilon/4$. All entropy, radius, and cardinality parameters are
polynomial in the stated quantities and independent of $N,t$.
Conditional on the pilot, the same finite masked-loss event works for
all widths and scales. Averaging over the pilot completes the proof.
\end{proof}

If the family has Vapnik--Chervonenkis (VC) dimension at most $h\ge1$, its restriction to the
finite Boolean cube has at most $(2^d+1)^h$ distinct subsets. The theorem
therefore applies with $\log M\le h(d+1)\log2$, even for a continuously
parameterized family. The family must be fixed before the validation and
labeled samples, or obtained from an independent pilot. No assertion is
made that arbitrary network instability sets have width-independent VC
dimension.

A concrete example is a union of neighborhoods of $M_0$ pilot centers,
with separate radii chosen by each network. Its uncovered region is an
upper orthant in the fixed vector of distances to the centers and has VC
dimension at most $M_0$. Thus network-dependent local coverage can be
handled with polynomial sample complexity. The prior still requires
actual stability throughout the covered regions. Pointwise testing of the
centers or observed validation inputs alone does not verify that condition.


\begin{corollary}[Using the training inputs to choose the support]
\label{cor:localized-same-inputs}
If the full exception family is fixed independently of $S$, the input
locations of $S$ itself may replace the independent validation sample in
\eqref{eq:localized-support}. The resulting input-dependent Gaussian prior,
followed by the ordinary training-margin likelihood, has the same
$\rho+\varepsilon$ guarantee when $n$ also satisfies the coverage sample
bound. No additional unlabeled sample is needed.
\end{corollary}
\begin{proof}
The coverage bound is uniform over the fixed family, and the masked-loss
bound is uniform over its product with the common finite candidate class.
Both events hold for the same iid input sample; their failure probabilities
can be added without independence. The component entropy and compression
arguments are deterministic statements for each realized sample.
\end{proof}

This corollary does not allow an arbitrarily complex exception family to be
learned from the same sample without accounting for its complexity. For
example, take as exception sets the complements of all unions of at most
$k$ Euclidean balls with centers in the Boolean cube. This family is fixed
in advance. On that cube each center
has at most $d+2$ distinct radius/inclusion choices, including an unused
ball, so the family has at most $[2^d(d+2)]^k$ descriptions. Posterior-chosen
centers and radii are therefore allowed at a cost
$\log M\le k[d\log2+\log(d+2)]$. This formalizes a condition that stability
holds on a few input neighborhoods covering most observed inputs, with a
statistical guarantee for their population coverage. It does not assert
that a small number of such neighborhoods suffices for any particular
high-dimensional data distribution.
